\par\indent\SourceLeaf{246}%
\textsc{Satz.} \textit{Zieht man auf einer gekrümmten Fläche von demselben Anfangspunkt aus unzählige kürzeste Linien gleicher Länge, so steht die Linie, die ihre Endpunkte verbindet, auf jeder von ihnen senkrecht.}

Wir hielten es für lohnend, diesen Satz aus der grundlegenden Eigenschaft der kürzesten Linien herzuleiten. Seine Richtigkeit lässt sich auch ohne Rechnung durch folgende Überlegung einsehen. Es seien \(AB,AB'\) zwei kürzeste Linien gleicher Länge, die bei \(A\) einen unendlich kleinen Winkel einschließen. Nehmen wir an, einer der Winkel, die das Element \(BB'\) mit den Linien \(BA,B'A\) bildet, weiche um einen endlichen Betrag von einem rechten Winkel ab. Nach dem Gesetz der Stetigkeit ist dann der eine Winkel größer als ein rechter Winkel, der andere kleiner. Der Winkel bei \(B\) sei \(=90^0-\omega\). Auf der Linie \(BA\) wählen wir einen Punkt \(C\) so, dass \(BC=BB'\mathord{.}\mathrm{cosec}\,\omega\) gilt. Da wir das unendlich kleine Dreieck \(BB'C\) als eben behandeln dürfen, gilt \(CB'=BC\mathord{.}\cos\omega\) und damit
\[
\begin{aligned}
AC+CB'&=AC+BC\mathord{.}\cos\omega=AB-BC\mathord{.}(1-\cos\omega)\\
&=AB'-BC(1-\cos\omega)
\end{aligned}
\]
das heißt, der Weg vom Punkt \(A\) zu \(B'\) über den Punkt \(C\) ist kürzer als die kürzeste Linie. Q. E. A.
\Needspace{10\baselineskip}\Article{16}
Dem Satz des vorigen Artikels fügen wir einen weiteren hinzu, den wir so formulieren. \textit{Denken wir uns auf einer gekrümmten Fläche eine beliebige Linie, von deren einzelnen Punkten unzählige kürzeste Linien gleicher Länge unter rechten Winkeln und nach derselben Seite ausgehen. Die Kurve, die ihre anderen Endpunkte verbindet, schneidet jede von ihnen unter einem rechten Winkel.} Zum Beweis genügt in der vorangegangenen Analyse eine einzige Änderung. Hier muss \(\varphi\) die von einem beliebigen Punkt aus gemessene Länge der \textit{gegebenen} Kurve bezeichnen oder, wenn Sie es vorziehen, eine Funktion dieser Länge. Die gesamte Argumentation gilt dann auch hier, mit der Änderung, dass \(S=0\) für \(r=0\) bereits in der Voraussetzung selbst enthalten ist. Dieser zweite Satz ist übrigens allgemeiner als der vorige, den man als darin enthalten ansehen kann, wenn man als gegebene Linie einen unendlich kleinen Kreis um den Mittelpunkt \(A\) wählt. Schließlich bemerken wir, dass auch hier geometrische Überlegungen an die Stelle der Analyse treten können. Da sie jedoch hinreichend naheliegen, gehen wir hier nicht näher darauf ein.
\Article{17}
Wir kehren zu der Formel \(\BareRoot{(E\mathrm{d}p^2+2F\mathrm{d}p\mathord{.}\mathrm{d}q+G\mathrm{d}q^2)}\) zurück, die allgemein
